\begin{proof}

\begin{enumerate}[label={$\langle1\rangle$\arabic*.},leftmargin=*]
\item The map $\ell(t)=(\log t)/(2\pi i)$ with values in $E$ is single-valued and
holomorphic on an annulus: changing the logarithm adds an integer.
Consequently $a_{\rm out}=a_{\rm in}+(\ell(t),0)$ is a holomorphic transition
for a principal $A$-bundle on $\Pp^1$. Product charts with compact fibre
prove that this bundle is Hausdorff and compact.

\item We choose disjoint discs about $0,1,\infty$ and replace their product models
by P4.1 using the holomorphic group isomorphisms
\[
 C_i\times B_i\longrightarrow A:\qquad
 (z,w)\mapsto(w,z),\quad(iz,iz+w),\quad(z,iz+w),
\]
respectively. Each lattice matrix $S_i$ displayed below has determinant
one. These maps, $\lambda_B=v\lambda_A$, the canonical punctured-disc
trivialization in P4.1, and the clutch $\ell$ specify all transitions.
Different base points separate in the base. Points over the same base
point lie in one Hausdorff local model, which separates them. The gluing is
therefore Hausdorff. Properness holds over each base chart and hence over
the base; compactness follows.

\item The two central components have connected elliptic fibres over connected
curves and meet each other. All remaining fibres equal $A$. The equations
$uv=t$ give exactly the indicated critical values. The punctured-disc
trivializations identify all fibre lattices with $\Lambda$, and the remaining
transition is a translation on $A$. Translation induces the identity on
$H_1(A;\Z)$, giving the asserted monodromy representation.

\end{enumerate}
\end{proof}
